% backmatter/source-map.tex

\addchap{Карта первоисточников}\label{ch:primary-sources}
Когда в~платёжной команде звучит фраза \enquote{надо проверить
в~правилах}, спор часто разгорается вокруг источника истины. Один человек
открывает свод правил сети, другой ищет ответ
в~технической спецификации, третий пересказывает заметку поставщика.
Одну и~ту же тему действительно могут описывать три класса
документов: публичные правила для участников, технические
спецификации и закрытые оперативные уведомления.

Отсюда первая задача навигации: прежде чем читать, понять, какой документ
в~этой теме главный, а~какой~--- только комментарий, пересказ или
вспомогательный портал интеграции.

\section*{EMVCo}\label{sec:ps-emvco}

EMVCo нужен там, где спор идёт не о~договорных обязанностях, а~о~поведении
платёжной техники. Если вопрос звучит как \enquote{что именно должны сделать
терминал, карта, токен или мобильный кошелёк}, искать ответ нужно прежде всего
здесь.

Если речь идёт о~контактной транзакции EMV, точкой входа остаются EMV Chip
Specifications, Books~1--4~\cite{emvco-specs} (см.~главу~\ref{ch:emv}). Для
бесконтактной оплаты нужен корпус EMV Contactless
Specifications~\cite{emvco-contactless-chip}, в~котором описаны NFC и~логика
бесконтактного ядра (см.~главу~\ref{ch:contactless}). Вопросы токенизации
закрывает EMV Payment Tokenisation
Specification~\cite{emvco-payment-tokenisation}: роли TSP, ограничения домена
токена и~его жизненный цикл (см.~главу~\ref{ch:tokenization}). Для 3DS основной
документ~--- EMV 3-D Secure Protocol and Core Functions
Specification~\cite{emvco-3ds}: сообщения AReq/ARes, сценарий дополнительной
проверки и~общая логика протокола (см.~главу~\ref{ch:3ds}).

EMVCo тяжелее всего читать подряд. Здесь выгоднее идти от вопроса к~разделу.
Если нужно понять жизненный цикл контактной операции, полезнее всего Book 3.
Если важно поведение терминала, кассира и интерфейса приёма, чаще открывают
Book 4. Если вопрос относится к~бесконтактному входу в~транзакцию, нужен
соответствующий том Contactless.

\section*{PCI SSC}\label{sec:ps-pci}

Стандарты PCI нужны, когда команда спорит о~периметре PCI и~наборе обязательных
контролей. Как ограничить себя SAQ~A, как контролировать JavaScript на странице
оплаты, как защищать PIN block и~где проводить точку дешифрования: на все эти
вопросы отвечать нужно документами PCI.

По периметру PCI и~базовому набору контролей нужен PCI~DSS~\cite{pci-dss-4}
(см.~главу~\ref{ch:pci-dss}). Для разных моделей интеграции существуют
сокращённые типы оценки SAQ~\cite{pci-saq}. Работу с~PIN и~HSM регулирует
PCI~PIN Security Requirements~\cite{pci-pin-security}. Когда вопрос упирается
в~шифрование от терминала до точки дешифрования, открывается
PCI~P2PE~\cite{pci-p2pe}.

Здесь практический порядок чтения один: сначала определить периметр PCI, потом
понять, какой путь оценки применим, и~затем разбирать отдельные требования.
Ошибка в~определении периметра обычно обходится дороже, чем неверное чтение
одного подпункта.

\section*{НСПК, Мир и~СБП}\label{sec:ps-nspk}

НСПК нужна там, где вопрос касается правил, операционной модели или формата
сообщения внутри РФ~--- для платёжной системы \enquote{Мир}, СБП или
внутрироссийского процессинга карт. Здесь полезно заранее определить, ищете вы
правовое поле, продуктовую модель или конкретный формат сообщения.

Правила платёжной системы~\enquote{Мир} задают роли и~обязанности
участников~\cite{nspk-mir-rules,nspk-mir-rules-v43}; на момент издания книги
действует редакция~4.3 от 3~декабря 2025 года. Вокруг них на портале для банков
\href{https://nspk.ru/banks/guides/}{\texttt{nspk.ru/banks/guides/}} лежит
набор документов: Стандарты ПС~\enquote{Мир} доопределяют требования к~системе
управления рисками информационной безопасности и~к~использованию карт
в~нефинансовых сервисах; тарифная политика, методические руководства
и~SAQ~D-MIR закрывают регуляторные и~операционные аспекты участия. Спецификация
бесконтактных платежей~\enquote{Мир} описывает приём.

Документы по СБП лежат на том же портале и~на
\href{https://sbp.nspk.ru/}{\texttt{sbp.nspk.ru}}. Правила оказания
операционных услуг и~услуг платёжного клиринга в~СБП задают работу
оператора~--- ОПКЦ~НСПК~\cite{sbp-main}; публичная часть OpenAPI ОПКЦ~СБП
выложена на
\href{https://sbp.nspk.ru/api}{\texttt{sbp.nspk.ru/api}}~\cite{nspk-sbp-openapi};
нормативная рамка подключения~--- документы Банка
России~\cite{cbr-sbp-actions}.

ОПКЦ~НСПК (операционный и~платёжный клиринговый центр)~--- функциональная роль
внутри НСПК: маршрутизация и~клиринг для~\enquote{Мир} и~для СБП, исторически
также внутрироссийская обработка операций международных платёжных систем.
Полные интерфейсные спецификации (форматы сообщений в~редакции НСПК, коды
расширений), процедуры подключения и~сертификационные материалы ОПКЦ
предоставляет участникам через закрытый портал поддержки. Программное
обеспечение НСПК (мобильные решения и~терминальные технологии) описано
в~открытом разделе~\cite{nspk-tsp}.

Если вопрос касается роли участника, продуктовой схемы или общей модели
работы~\enquote{Мир} либо СБП, обычно хватает публичных правил и~стандартов.
Если же спор упирается в~формат сообщения, сертификацию или детали
внутрироссийского процессинга, понадобится закрытая документация ОПКЦ через
участника системы или процессингового партнёра.

\section*{Банк России}\label{sec:ps-cbr}

Документы Банка России нужны, когда требуется понять рамки допустимого
поведения. Если вопрос звучит как \enquote{обязаны ли мы}, \enquote{имеем ли мы
право} или \enquote{какой статус нужен}, переходить следует к~нормативным актам
ЦБ и~федеральным законам.

Основной публичный раздел~---
\href{https://www.cbr.ru/PSystem/}{\texttt{cbr.ru/PSystem/}}~--- собирает
вместе нормативные акты, методические рекомендации, информационные письма,
реестры и~справочные материалы по национальной платёжной системе: разделы
\enquote{Платёжная система Банка России}, \enquote{Основные документы},
\enquote{Стандарт ISO~20022}, \enquote{Система передачи финансовых сообщений},
\enquote{Система быстрых платежей}, \enquote{Надзор и~наблюдение},
\enquote{Доступ оператора платёжной системы}~\cite{cbr-ps}. Проекты нормативных
актов выходят в~отдельном разделе публичного обсуждения
\href{https://www.cbr.ru/project_na/}{\texttt{cbr.ru/project\_na/}}
с~указанными сроками для замечаний до принятия документа~\cite{cbr-project-na}.

161-ФЗ задаёт базовые роли платёжного рынка~\cite{fz-161}. 762-П описывает
правила перевода денежных средств~\cite{cbr-762p}. 821-П задаёт защиту
информации при~переводах денежных средств, 850-П~--- операционную надёжность
кредитных организаций, 851-П~--- противодействие переводам без~согласия
клиента; все три отсылают к~57580.1 в~части состава мер и~к~57580.2 в~части
оценки соответствия~\cite{cbr-821p,cbr-850p,cbr-851p}. 266-П удерживает
карточную специфику~\cite{cbr-266p}; на смену готовится новый акт Банка
России~--- проект Положения \enquote{О порядке предоставления электронных
средств платежа и~осуществления операций с~их использованием} (последняя
  публичная редакция~--- от 28~октября 2024~года; на момент написания этих строк
в~силу не вступил). Для цифрового рубля отправная точка~---
340-ФЗ~\cite{fz-340}; требования к~защите информации для участников платформы
установлены 833-П~\cite{cbr-833p}. Для СБП и~других специальных тем набор
документов свой.

Помимо нормативных актов, ЦБ выпускает методические рекомендации (серии~МР)
и~информационные письма~--- они не задают обязанности, но фиксируют позицию
регулятора по спорным вопросам. В~ежедневной работе их полезно отслеживать
вместе с~основными актами: интерпретация часто появляется раньше, чем
формальное изменение акта.

Нормативные акты полезно читать связкой: закон задаёт роли и~пределы модели,
положение описывает операционную логику, письма и разъяснения показывают
регуляторскую интерпретацию спорных мест. Для практической работы удобнее
пользоваться правовыми системами, где доступны история изменений и перекрёстные
ссылки.

\section*{ISO}\label{sec:ps-iso}

ISO нужен, когда требуется обобщённый язык, более широкий, чем диалект
конкретной сети. Прикладная задача~--- сначала сетевой диалект, потом ISO как
верхний уровень; архитектурная или межсистемная (новый обмен)~--- сразу
ISO~20022 как исходную модель.

ISO~8583 остаётся базовым стандартом структуры финансовых
сообщений~\cite{iso-8583-2023}; на практике его конкретизируют сетевые
реализации. ISO~20022~\cite{iso-20022} задаёт современную модель финансовых
сообщений и~особенно важен для международных систем, SWIFT~MX и~многих систем
мгновенных платежей (см.~раздел~\ref{sec:mc-correspondent}).
ISO~4217~\cite{iso-4217}, ISO~3166--1~\cite{iso-3166-1}
и~ISO~9564--1~\cite{iso-9564} держат базовые коды валют, коды стран и~форматы
PIN~block.
